% =============================================================================
% 视觉应用软件开发技术实践——指导说明书【最终全优化排版版】
% 优化点：
% 1. 删除封面多余文档整合说明，封面留白重新均衡分配
% 2. 全局排版统一、层级间距梯度标准化、段落缩进规范
% 3. 代码块限定宽度防横向溢出，自动换行柔和
% 4. 表格宽度压缩、行高适中，彻底消除 Overfull \hbox
% 5. 目录增加引导点线，页码对齐美观
% 6. 页眉无报错，弱化xeCJK字体重定义警告
% 7. 列表紧凑无冗余空白，正文段落间距舒适
% 8. 所有下划线转义，无Missing $编译错误
% 9. 移除非法CJKfontfamily，无未定义控制序列报错
% 10. 页面边距、章节、代码、图表留白分层清晰
% 编译方式：XeLaTeX 仅存在VS Code cSpell无害拼写提示，无LaTeX致命报错
% =============================================================================
\documentclass[UTF8,a4paper,12pt]{ctexart}

% -------------------------- 全局页面基础设置 --------------------------
\usepackage[top=2.5cm, bottom=2.5cm, left=2.8cm, right=2.8cm]{geometry}
\usepackage{indentfirst}               % 中文段落首行自动缩进
\setlength{\parindent}{2em}            % 首行缩进2个汉字
\setlength{\parskip}{0.3em}            % 段间微小留白，阅读更舒适

% -------------------------- 目录美化（带引导点线） --------------------------
\usepackage{tocloft}
\renewcommand{\cftsecfont}{\bfseries}
\renewcommand{\cftsecpagefont}{\bfseries}
\renewcommand{\cftsecleader}{\cftdotfill{\cftdotsep}}
\renewcommand{\cftsubsecleader}{\cftdotfill{\cftdotsep}}
\setlength{\cftbeforesecskip}{8pt}
\setlength{\cftbeforesubsecskip}{4pt}
\setlength{\cfttabindent}{0pt}

% -------------------------- 字体配置（弱化xeCJK重定义警告） --------------------------
\usepackage{fontspec}
\setCJKmainfont{SimSun}[AutoFakeBold=2.5]
\setmainfont{Times New Roman}

% -------------------------- 代码高亮包 + 图表标题格式 --------------------------
\usepackage{listings}
\usepackage{xcolor}
\usepackage{caption}
\captionsetup{
    tableposition=top,
    figureposition=bottom,
    font=small,
    labelfont=bf,
    skip=8pt
}

% 代码配色定义
\definecolor{codebg}{RGB}{245,245,245}
\definecolor{codeframe}{RGB}{180,180,180}
\definecolor{codekeyword}{RGB}{0,0,180}
\definecolor{codecomment}{RGB}{0,130,0}
\definecolor{codestring}{RGB}{160,20,20}

% HALCON专用代码样式（限制宽度防止页面溢出）
\lstdefinestyle{halconstyle}{
    backgroundcolor=\color{codebg},
    frame=single,
    rulecolor=\color{codeframe},
    basicstyle=\ttfamily\small,
    keywordstyle=\color{codekeyword}\bfseries,
    commentstyle=\color{codecomment}\itshape,
    stringstyle=\color{codestring},
    tabsize=4,
    showstringspaces=false,
    breaklines=true,
    breakatwhitespace=true,
    numbers=left,
    numberstyle=\tiny\color{gray},
    numbersep=6pt,
    captionpos=b,
    aboveskip=10pt,
    belowskip=10pt,
    language=C++,
    linewidth=0.92\textwidth,
    morekeywords={
        read_image,dev_close_window,dev_open_window_fit_image,dev_set_window,
        dev_display,rgb3_to_gray,draw_rectangle1,gen_rectangle1,reduce_domain,
        threshold,opening_circle,connection,select_shape,shape_trans,
        smallest_rectangle2,count_obj,sort_region,gen_rectangle2,concat,
        disp_message,dev_update_off,dev_set_draw,dev_set_color,
        dev_set_line_width,set_display_font,rgb1_to_gray,
        create_data_code_2d_model,find_data_code_2d,dev_open_window,
        for,endif,endfor,if,elseif
    }
}

% -------------------------- 表格、图片自适应 --------------------------
\usepackage{graphicx}
\usepackage{float}
\usepackage{tabularx}
\usepackage{array}
\renewcommand{\arraystretch}{1.15}    % 表格行高适中，不拥挤不松散

% -------------------------- 超链接（无边框） --------------------------
\usepackage[hidelinks]{hyperref}

% -------------------------- 有序/无序列表统一紧凑格式 --------------------------
\usepackage{enumitem}
\setlist{nosep,leftmargin=2em,itemsep=2pt,topsep=3pt}

% -------------------------- 页眉页脚（彻底消除headheight警告） --------------------------
\usepackage{fancyhdr}
\pagestyle{fancy}
\setlength{\headheight}{14pt}
\fancyhf{}
\fancyhead[L]{\small 视觉应用软件开发技术实践——指导说明书}
\fancyhead[R]{\small 工程训练（金工实习）选修模块}
\fancyfoot[C]{\thepage}
\renewcommand{\headrulewidth}{0.5pt}
\renewcommand{\footrulewidth}{0pt}

% -------------------------- 章节标题间距分层控制 --------------------------
\ctexset{
    section={
        format={\Large\bfseries},
        beforeskip=24pt,
        afterskip=12pt,
    },
    subsection={
        format={\large\bfseries},
        beforeskip=16pt,
        afterskip=8pt,
    },
    subsubsection={
        format={\normalsize\bfseries},
        beforeskip=12pt,
        afterskip=6pt,
    },
    paragraph={
        format={\normalsize\bfseries},
        beforeskip=8pt,
        afterskip=4pt,
    }
}

\begin{document}

% ============================== 封面（已删除多余文档说明，留白均衡） ==============================
\begin{titlepage}
\thispagestyle{empty}
\begin{center}
\vspace*{5cm}

{\Huge\bfseries 视觉应用软件开发技术实践}
\vspace{1cm}
{\LARGE 指导说明书}

\vspace{2.8cm}
{\large 工程训练（金工实习）——选修模块}
\vspace{0.6cm}
{\large 小班实训，单次最多容纳10人}

\vfill
\vspace{3cm}
\end{center}
\end{titlepage}

% ============================== 目录页 ==============================
\tableofcontents
\newpage
\setcounter{page}{1}

% ===================================================================
\section{开发环境配置与图像基础操作}
\subsection{实验内容与目的}
\begin{itemize}
    \item 掌握 HALCON、Visual Studio C++ 完整开发环境搭建流程；
    \item 实现图像读取、窗口显示、像素空间基础操作；
    \item 理解数字图像存储原理与RGB/灰度色彩空间转换逻辑。
\end{itemize}

\subsection{实验设备清单}
\begin{itemize}
    \item \textbf{硬件设备}：机器视觉综合教学平台、配套工业相机；
    \item \textbf{软件工具}：MVTec HALCON、Visual Studio 2022（C++控制台开发）。
\end{itemize}

% ===================================================================
\section{实验任务一：Blob分析实现PCB焊片区域提取}
\subsection{任务概述}
使用工业相机拍摄标准焊板样本图像，拍摄要求：焊板水平摆放、无大幅倾斜、画面清晰无过曝/欠曝。完整检测流程采用标准Blob图像分析流水线，分步完成区域提取、形态学去噪、连通域筛选、缺陷判定。

\subsection{分步操作流程}
\subsubsection{步骤1：图像读入与自适应窗口显示}
打开HALCON程序编辑器，调用\texttt{read\_image}算子加载本地图像；支持单步逐算子调试。读取图像尺寸后自适应创建图形窗口并绑定显示。
\begin{lstlisting}[style=halconstyle, caption={图像读入与自适应窗口显示}]
* 读取本地样本图像，自适应创建窗口并展示
read_image(Image, 'E:/实验教学/2025工训选修模块/焊接片/0819/barcode12.jpg')
* 关闭软件默认空白图形窗口
dev_close_window()
* 根据图像分辨率自动适配窗口尺寸
dev_open_window_fit_image(Image, 0, 0, -1, -1, WindowHandle)
* 绑定当前操作窗口句柄
dev_set_window(WindowHandle)
* 在窗口中渲染原图
dev_display(Image)
\end{lstlisting}

\subsubsection{步骤2：灰度转换与自定义ROI裁剪}
将三通道彩色图像转为单通道灰度图；交互式手绘矩形划定待检测焊片区域，生成ROI并对原图裁剪，仅保留目标检测区域，再基于灰度直方图完成阈值分割。
\begin{lstlisting}[style=halconstyle, caption={彩色图灰度化 + ROI感兴趣区域提取}]
* RGB三通道图像转换单通道灰度图
rgb3_to_gray(Image, Image, Image, ImageGray)
dev_display(ImageGray)
* 交互式鼠标手绘矩形检测框
draw_rectangle1(WindowHandle, Row1, Column1, Row2, Column2)
* 根据手绘坐标生成矩形ROI区域对象
gen_rectangle1(ROI_0, Row1, Column1, Row2, Column2)
* 依据ROI裁剪灰度图像，减少无效计算区域
reduce_domain(ImageGray, ROI_0, ImageReduced)
* 灰度阈值分割，提取高亮焊片区域
threshold(ImageReduced, Regions, 190, 255)
\end{lstlisting}

\paragraph{灰度直方图工具使用规范}
执行\texttt{reduce\_domain}完成图像裁剪后，在变量窗口右键\texttt{ImageReduced}，选择「清除并显示变量」；点击工具栏「灰度直方图」图标，绿色标线为灰度下限、红色标线为灰度上限，拖动两条标线框选目标灰度区间后，点击「阈值」激活分割预览，最后点击「插入代码」自动生成对应threshold算子。

\subsubsection{步骤3：形态学去噪与连通域特征筛选}
开运算用于消除微小噪点、划痕，保留大面积焊片轮廓；拆分独立连通域后按面积特征过滤干扰区域，再转换为最小外接矩形用于尺寸、倾斜度检测。
\begin{lstlisting}[style=halconstyle, caption={形态学开运算 + 连通域特征筛选}]
* 圆形结构元素开运算，去除微小噪点、细碎干扰
opening_circle(Regions, RegionOpening, 3.5)
* 拆分互不连通的独立区域
connection(RegionOpening, ConnectedRegions)
* 按面积特征筛选有效焊片区域，过滤过小杂点
select_shape(ConnectedRegions, SelectedRegions, 'area', 'and', 2572.73, 5000)
* 将不规则区域统一转换为标准矩形
shape_trans(SelectedRegions, RegionTrans, 'rectangle2')
* 计算矩形中心坐标、倾斜角度、长宽尺寸
smallest_rectangle2(RegionTrans, Row, Column, Phi, Length1, Length2)
\end{lstlisting}

\paragraph{特征直方图筛选操作}
连通域拆分完成后，右键连通域变量选择「清除并显示」；打开工具栏「特征直方图」，横轴选择面积特征，拖动滑块限定有效面积区间，点击「插入代码」自动生成\texttt{select\_shape}筛选逻辑。

\subsubsection{步骤4：焊片缺陷判定逻辑}
统计焊片总数、遍历判断单块焊片倾斜角度；总数不足10判定为缺焊，单块倾斜弧度大于0.08判定为偏斜不良，分别用红色框标注缺陷并输出文字提示。
\begin{lstlisting}[style=halconstyle, caption={焊片数量、倾斜缺陷判定完整逻辑}]
* 重置绘图样式：仅绘制区域轮廓
dev_display(Image)
dev_set_draw('margin')
dev_set_color('green')
dev_set_line_width(3)
set_display_font(WindowHandle, 24, 'mono', 'true', 'false')

count_obj(RegionTrans, Number)
sort_region(RegionTrans, SortedRegions, 'first_point', 'true', 'row')
flag := 1

* 循环遍历所有焊片，检测倾斜角度
for Index := 1 to Number by 1
    if (abs(Phi[Index-1]) > 0.08)
        gen_rectangle2(Rectangle1, Row[Index-1], Column[Index-1],
                       Phi[Index-1], Length1[Index-1], Length2[Index-1])
        dev_set_color('red')
        dev_display(Rectangle1)
        disp_message(WindowHandle, Phi[Index-1], 'image',
                     Row[Index-1], Column[Index-1], 'black', 'true')
        flag := 0
    endif
endfor

* 分情况判定输出检测结果
if (Number < 10)
    disp_message(WindowHandle, 'Quantity error', 'image',
                 Row[0], Column[0], 'red', 'true')
    flag := 0
elseif (flag == 1)
    disp_message(WindowHandle, 'OK', 'image',
                 Row[0], Column[0], 'blue', 'true')
endif

dev_update_off()
\end{lstlisting}

\subsection{拓展功能：二维码学生信息识别}
调用2D二维码解码算子读取样本附带学生信息QR码，在窗口左上角打印解析出的学号；使用HALCON标准\texttt{concat}拼接文本。
\begin{lstlisting}[style=halconstyle, caption={QR二维码识别与学号显示}]
* 读取带二维码的样本图像
read_image(Image, 'C:/Users/lenovo/Desktop/2024312011.jpg')
* 新建独立窗口展示二维码画面
dev_open_window_fit_image(Image, 0, 0, -1, -1, WindowHandle1)
dev_set_window(WindowHandle1)
dev_display(Image)
* 彩色图转灰度图用于解码
rgb1_to_gray(Image, GrayImage)
* 初始化QR Code二维码识别模型
create_data_code_2d_model('QR Code', [], [], DataCodeHandle)
* 灰度图像中查找并解码二维码
find_data_code_2d(GrayImage, SymbolXLDs, DataCodeHandle,
                  [], [], ResultHandles, DecodedDataStrings)
set_display_font(WindowHandle1, 24, 'mono', 'true', 'false')
concat('学生二维码：', DecodedDataStrings, ShowText)
disp_message(WindowHandle1, ShowText, 'window', 12, 12, 'black', 'true')
\end{lstlisting}

\subsection{实验任务一核心HALCON算子汇总}
\begin{table}[H]
\centering
\caption{任务一核心HALCON算子功能对照表}
\begin{tabularx}{0.94\textwidth}{|l|X|}
\hline
\textbf{算子名称} & \textbf{功能详细说明} \\
\hline
\texttt{read\_image} & 从本地磁盘读取图像文件至内存变量 \\
\texttt{dev\_close\_window} & 关闭软件启动自带默认图形窗口 \\
\texttt{dev\_open\_window\_fit\_image} & 根据图像分辨率自适应创建可视化窗口 \\
\texttt{dev\_set\_window} & 绑定窗口句柄，指定后续绘图操作目标窗口 \\
\texttt{dev\_display} & 在绑定窗口渲染图像、区域、线条等图形对象 \\
\texttt{rgb3\_to\_gray} & 三通道RGB彩色图像转换单通道灰度图像 \\
\texttt{draw\_rectangle1} & 鼠标交互式手绘矩形，输出对角坐标 \\
\texttt{gen\_rectangle1} & 根据对角坐标生成矩形ROI区域对象 \\
\texttt{reduce\_domain} & 按ROI区域裁剪原图，仅保留指定检测范围 \\
\texttt{threshold} & 灰度阈值分割，分离前景目标与背景 \\
\texttt{opening\_circle} & 圆形结构元素形态学开运算，去除微小噪点 \\
\texttt{connection} & 拆分图像中互不连通的独立区域对象 \\
\texttt{select\_shape} & 基于面积、长宽等几何特征过滤区域 \\
\texttt{shape\_trans} & 将不规则区域统一转换为矩形、椭圆等标准轮廓 \\
\texttt{smallest\_rectangle2} & 计算区域最小外接矩形，输出中心、角度、长宽 \\
\texttt{count\_obj} & 统计区域、XLD等图形对象总数量 \\
\texttt{sort\_region} & 按坐标、面积等规则对区域集合排序 \\
\texttt{gen\_rectangle2} & 根据中心、倾斜角、长宽生成任意角度矩形 \\
\texttt{disp\_message} & 在图形窗口指定位置打印文本提示信息 \\
\texttt{dev\_set\_draw} & 设置区域绘制模式：轮廓/填充 \\
\texttt{dev\_set\_color} & 修改图形绘制线条、填充的显示颜色 \\
\texttt{dev\_set\_line\_width} & 调整轮廓绘制线条像素宽度 \\
\texttt{create\_data\_code\_2d\_model} & 初始化2D二维码解码识别模型 \\
\texttt{find\_data\_code\_2d} & 在灰度图像中定位并解析二维码内容 \\
\hline
\end{tabularx}
\end{table}

% ===================================================================
\newpage
\section{实验任务二：VS2022开发独立视觉应用软件}
\subsection{任务概述}
HALCON离线调试代码验证效果无误后，导出标准C++源码，在Visual Studio 2022中配置HALCON依赖环境，封装编译生成独立可执行视觉检测程序，脱离HALCON软件直接运行。所有图像业务逻辑统一封装在导出代码\texttt{action()}函数内。

\subsection{步骤1：HALCON工程导出标准C++代码}
HALCON完整流程调试无报错后，通过软件顶部菜单栏导出功能，将HDevelop脚本导出为独立\texttt{.cpp}源文件；图像读取、形态学、缺陷判定全部封装在\texttt{action()}函数内部，是VS开发的核心逻辑载体。

\subsection{步骤2：VS控制台项目创建与依赖环境配置}
\subsubsection{新建空白C++控制台项目}
\begin{enumerate}
    \item 启动Visual Studio 2022，首页点击「创建新项目」；
    \item 模板筛选选择「控制台应用（C++）」，点击下一步；
    \item 自定义项目名称（例：\texttt{VisionPCBInspector}），选择本地存放路径，完成创建。
\end{enumerate}

\subsubsection{配置HALCON头文件包含目录}
右键左侧解决方案资源管理器项目名称，选择「属性」打开配置页：
\begin{itemize}
    \item 展开 \textbf{C/C++} → \textbf{常规}；
    \item 编辑「附加包含目录」，填入本地HALCON安装下include文件夹路径。
\end{itemize}

\subsubsection{配置HALCON静态库目录}
\begin{itemize}
    \item 展开 \textbf{链接器} → \textbf{常规}；
    \item 编辑「附加库目录」，填入对应x64平台lib库文件夹路径。
\end{itemize}

\subsubsection{添加链接依赖库}
\begin{itemize}
    \item 展开 \textbf{链接器} → \textbf{输入}；
    \item 在「附加依赖项」末尾追加 \texttt{halconcpp.lib}。
\end{itemize}

\subsection{步骤3：源码导入与项目编译生成程序}
环境配置全部完成后执行如下操作：
\begin{enumerate}
    \item 打开HALCON导出的\texttt{.cpp}文件，全选复制所有代码，覆盖VS项目默认主cpp文件；
    \item 顶部菜单栏选择「生成」→「重新生成解决方案」，执行完整编译；
    \item 无语法、依赖报错时，在项目输出目录生成exe可执行程序。
\end{enumerate}

\subsection{步骤4：程序调试运行与闪退问题解决方案}
\begin{itemize}
    \item 菜单栏点击「调试」→「开始执行（不调试）」，启动视觉检测程序查看画面输出。
\end{itemize}

\paragraph{高频故障：程序运行后窗口直接闪退}
故障原因：原生导出代码无阻塞等待逻辑，图像渲染完成后主线程立即结束进程。
推荐两种修复方案，任选其一添加至主函数末尾：
\begin{itemize}
    \item 方案1：延时阻塞 \texttt{Sleep(2000)}；窗口保持2秒后自动关闭；
    \item 方案2：键盘阻塞 \texttt{getchar()}；按下任意键才关闭图像窗口，便于观察检测结果。
\end{itemize}

\subsection{VS2022 HALCON环境配置速查表}
\begin{table}[H]
\centering
\caption{VS2022 HALCON依赖一键配置清单}
\begin{tabularx}{0.94\textwidth}{|l|l|X|}
\hline
\textbf{配置分类} & \textbf{属性层级路径} & \textbf{填写内容示例} \\
\hline
附加包含目录 & C/C++ → 常规 & \texttt{F:\textbackslash Program Files\textbackslash MVTec\textbackslash HALCON\textbackslash include} \\
附加库目录 & 链接器 → 常规 & \texttt{F:\textbackslash Program Files\textbackslash MVTec\textbackslash HALCON\textbackslash lib\textbackslash x64-win64} \\
附加依赖项 & 链接器 → 输入 & 末尾追加 \texttt{halconcpp.lib} \\
\hline
\end{tabularx}
\end{table}

% ===================================================================
\newpage
\section{附录：PCB焊片缺陷检测完整HALCON参考源码}
以下为不含二维码拓展功能的完整Blob检测代码，供实训学生对照调试：
\begin{lstlisting}[style=halconstyle, caption={PCB LED焊片缺陷检测完整参考代码}]
* ========== 1. 图像读取与自适应窗口 ==========
read_image(Image, 'E:/实验教学/2025工训选修模块/焊接片/0819/barcode12.jpg')
dev_close_window()
dev_open_window_fit_image(Image, 0, 0, -1, -1, WindowHandle)
dev_set_window(WindowHandle)
dev_display(Image)

* ========== 2. RGB图像转灰度 ==========
rgb3_to_gray(Image, Image, Image, ImageGray)
dev_display(ImageGray)

* ========== 3. 交互式ROI区域裁剪 ==========
draw_rectangle1(WindowHandle, Row1, Column1, Row2, Column2)
gen_rectangle1(ROI_0, Row1, Column1, Row2, Column2)
reduce_domain(ImageGray, ROI_0, ImageReduced)

* ========== 4. 灰度阈值分割提取焊片前景 ==========
threshold(ImageReduced, Regions, 190, 255)

* ========== 5. 形态学开运算去除微小噪点 ==========
opening_circle(Regions, RegionOpening, 3.5)

* ========== 6. 连通域拆分 ==========
connection(RegionOpening, ConnectedRegions)

* ========== 7. 按面积筛选有效焊片区域 ==========
select_shape(ConnectedRegions, SelectedRegions,
             'area', 'and', 2572.73, 5000)

* ========== 8. 区域转换最小外接矩形 ==========
shape_trans(SelectedRegions, RegionTrans, 'rectangle2')
smallest_rectangle2(RegionTrans, Row, Column, Phi, Length1, Length2)

* ========== 9. 缺陷可视化判定与绘制 ==========
dev_display(Image)
dev_set_draw('margin')
dev_set_color('green')
dev_set_line_width(3)
set_display_font(WindowHandle, 24, 'mono', 'true', 'false')

count_obj(RegionTrans, Number)
sort_region(RegionTrans, SortedRegions, 'first_point', 'true', 'row')
flag := 1

for Index := 1 to Number by 1
    if (abs(Phi[Index-1]) > 0.08)
        gen_rectangle2(Rectangle1, Row[Index-1], Column[Index-1],
                       Phi[Index-1], Length1[Index-1], Length2[Index-1])
        dev_set_color('red')
        dev_display(Rectangle1)
        disp_message(WindowHandle, Phi[Index-1], 'image',
                     Row[Index-1], Column[Index-1], 'black', 'true')
        flag := 0
    endif
endfor

if (Number < 10)
    disp_message(WindowHandle, 'Quantity error', 'image',
                 Row[0], Column[0], 'red', 'true')
elseif (flag == 1)
    disp_message(WindowHandle, 'OK', 'image',
                 Row[0], Column[0], 'blue', 'true')
endif

dev_update_off()
\end{lstlisting}

\vspace{3cm}
\begin{center}
\rule{0.6\textwidth}{0.5pt}\\[8pt]
{\small 工程训练（金工实习）选修模块 · 视觉应用软件开发技术实践指导说明书}
\end{center}

\end{document}